% !TeX root = ../main.tex

\begin{figure}[htbp]
	\centering
	\begin{tikzpicture}[
		font=\footnotesize,
		ag/.style={
			rectangle,
			rounded corners=2pt,
			draw,
			thick,
			minimum width=39mm,
			minimum height=10mm,
			align=center
		},
		ciego/.style={
			ag,
			very thick,
			fill=black!8
		},
		estado/.style={
			rectangle,
			draw,
			dashed,
			minimum width=120mm,
			minimum height=7mm,
			align=center,
			fill=black!3
		},
		nota/.style={
			rectangle,
			draw,
			dotted,
			minimum width=120mm,
			minimum height=7mm,
			align=center
		},
		fl/.style={
			-{Stealth[length=1.8mm]}
		}
		]
		
		% Génesis
		\node[ciego] (g1) at (0,0)
		{\texttt{agent\_1} (ciego)\\
			$\text{cesta}\rightarrow r_1^{(0)}$};
		
		\node[ag] (g2) at (4.3,0)
		{\texttt{agent\_2}\\
			$\text{cesta}\rightarrow r_2^{(0)}$};
		
		\node[ag] (g3) at (8.6,0)
		{\texttt{agent\_3}\\
			$\text{cesta}\rightarrow r_3^{(0)}$};
		
		\node[left=3mm of g1, align=right]
		{$t=0$\\génesis};
		
		\node[estado] (e0) at (4.3,-1.5)
		{Estado $t=0$:
			$\{r_1^{(0)},r_2^{(0)},r_3^{(0)}\}$};
		
		\draw[fl] (g1) -- (e0);
		\draw[fl] (g2) -- (e0);
		\draw[fl] (g3) -- (e0);
		
		% Primer turno
		\node[ciego] (a1) at (0,-3.2)
		{\texttt{agent\_1} (ciego)\\
			$\text{cesta} + r_2^{(0)} + r_3^{(0)}
			\rightarrow r_1^{(1)}$};
		
		\node[ag] (a2) at (4.3,-3.2)
		{\texttt{agent\_2}\\
			$\text{cesta} + r_1^{(0)} + r_3^{(0)}
			\rightarrow r_2^{(1)}$};
		
		\node[ag] (a3) at (8.6,-3.2)
		{\texttt{agent\_3}\\
			$\text{cesta} + r_1^{(0)} + r_2^{(0)}
			\rightarrow r_3^{(1)}$};
		
		\node[left=3mm of a1, align=right]
		{$t=1$};
		
		\draw[fl] (e0) -- (a1);
		\draw[fl] (e0) -- (a2);
		\draw[fl] (e0) -- (a3);
		
		\node[estado] (e1) at (4.3,-4.7)
		{Estado $t=1$:
			$\{r_1^{(1)},r_2^{(1)},r_3^{(1)}\}$};
		
		\draw[fl] (a1) -- (e1);
		\draw[fl] (a2) -- (e1);
		\draw[fl] (a3) -- (e1);
		
		\node[nota] (continuacion) at (4.3,-5.9)
		{$t=2,3,4$: se repite el esquema utilizando
			únicamente las respuestas ajenas de $t-1$.};
		
		\draw[fl] (e1) -- (continuacion);
		
	\end{tikzpicture}
	
	\caption[Protocolo de deliberación]{
		Secuencia de interacción del comité, donde $r_i^{(t)}$
		representa la respuesta del agente $i$ en el turno $t$.
		En génesis, los agentes deciden de forma independiente.
		En los turnos posteriores, cada agente recibe la cesta y
		las respuestas de los otros dos agentes en el turno anterior,
		pero no su propia respuesta ni las de rondas previas.
		El estado se actualiza después de que los tres hayan respondido,
		por lo que la actualización es síncrona.
	}
	\label{fig:protocolo_deliberacion}
\end{figure}